% corpus_supplement: Erdős problem #114
% source: https://www.erdosproblems.com/latex/114
% fetched_at: 2026-07-14T18:48:30Z
% payload_sha256: ab97c460fae09f7b1d261b5d6c150f9461118a8f5241687d8bd6a053b6637e6a
% statement/remarks/references extracted by erdos_ingest.extract_source_page;
% rendered by erdos_ingest.render_tex — do not edit
\documentclass{article}
\usepackage{amsmath, amssymb, amsthm}
\usepackage{hyperref}
\usepackage{geometry}
\geometry{margin=1in}

\title{Erd\H{o}s Problem \#114}
\date{Last updated: 2025-12-28}
\author{Source: erdosproblems.com}

\begin{document}
\maketitle

\noindent\textbf{Status:} OPEN \quad \textbf{No prize}

\noindent\textbf{Tags:} polynomials, analysis

\medskip
\noindent\textbf{Problem Statement:}

If $p(z)\in\mathbb{C}[z]$ is a monic polynomial of degree $n$ then is the length of the curve $\{ z\in \mathbb{C} : \lvert p(z)\rvert=1\}$ maximised when $p(z)=z^n-1$?

\bigskip
\noindent\textbf{Remarks:}

A problem of Erd\H{o}s, Herzog, and Piranian \cite{EHP58}. It is also listed as Problem 4.10 in \cite{Ha74}, where it is attributed to Erd\H{o}s.

In \cite{Va99} it is just aked whether the length is at most $2n+O(1)$. This is true, as a consequence of result of Tao \cite{Ta25} below.

Let the maximal length of such a curve be denoted by $f(n)$.
{UL}
{LI}The length of the curve when $p(z)=z^n-1$ is $2n+O(1)$, and hence the conjecture implies in particular that $f(n)=2n+O(1)$.{/LI}
{LI}Dolzhenko \cite{Do61} proved $f(n) \leq 4\pi n$, but few were aware of this work.{/LI}
{LI}Pommerenke \cite{Po61} proved $f(n)\ll n^2$.{/LI}
{LI}Borwein \cite{Bo95} proved $f(n)\ll n$ (Borwein was unaware of Dolzhenko's earlier work). The prize of \$250 is reported by Borwein \cite{Bo95}.{/LI}
{LI}Eremenko and Hayman \cite{ErHa99} proved the full conjecture when $n=2$, and $f(n)\leq 9.173n$ for all $n$.{/LI}
{LI}Danchenko \cite{Da07} proved $f(n)\leq 2\pi n$.{/LI}
{LI}Fryntov and Nazarov \cite{FrNa09} proved that $z^n-1$ is a local maximiser, and solved this problem asymptotically, proving that\[f(n)\leq 2n+O(n^{7/8}).\]{/LI}
{LI} Tao \cite{Ta25} has proved that $p(z)=z^n-1$ is the unique (up to rotation and translation) maximiser for all sufficiently large $n$.
{/UL}
Erd\H{o}s, Herzog, and Piranian \cite{EHP58} also ask whether the length is at least $2\pi$ if $\{ z: \lvert f(z)\rvert<1\}$ is connected (which $z^n$ shows is the best possible). This was proved by Pommerenke \cite{Po59}.

\bigskip
\noindent\textbf{References:}

[Bo95] Borwein, Peter, The arc length of the lemniscate {$\{|p(z)|=1\}$}. Proc. Amer. Math. Soc. (1995), 797--799.
 
 [Da07] Danchenko, V. I., The lengths of lemniscates. {V}ariations of rational
functions. Mat. Sb. (2007), 51--58.
 
 [Do61] Dol\v zenko, E. P., Some estimates concerning algebraic hypersurfaces and
derivatives of rational functions. Dokl. Akad. Nauk SSSR (1961), 1287--1290.
 
 [EHP58] Erd\H{o}s, P. and Herzog, F. and Piranian, G., Metric properties of polynomials. J. Analyse Math. (1958), 125-148.
 
 [ErHa99] Eremenko, Alexandre and Hayman, Walter, On the length of lemniscates. Michigan Math. J. (1999), 409--415.
 
 [FrNa09] Fryntov, Alexander and Nazarov, Fedor, New estimates for the length of the {E}rd\H
os-{H}erzog-{P}iranian lemniscate. (2009), 49--60.
 
 [Ha74] Hayman, W. K., Research problems in function theory: new problems. (1974), 155--180.
 
 [Po59] Pommerenke, Ch., On some problems by Erd\H{o}s, Herzog and Piranian. Michigan Math. J. (1959), 221-225.
 
 [Po61] Pommerenke, Ch., On metric properties of complex polynomials. Michigan Math. J. (1961), 97-115.
 
 [Ta25] T. Tao, The maximal length of the Erd\H{o}s-Herzog-Piranian leminscate length in high degree. arXiv:2512.12455 (2025).
 
 [Va99] Various, Some of Paul's favorite problems. Booklet produced for the conference "Paul Erd\H{o}s and his mathematics", Budapest, July 1999 (1999).

\bigskip
\noindent\small{Source: \url{https://www.erdosproblems.com/114}}
\end{document}
